\begin{textsample}{POS Dim 3 – se – Score 28.00 – se/se\_inf\_8.txt}  \label{ex:f3_pos_162}
2.4 Elementos Para a Organização Curricular 2.4.1 Direitos de Aprendizagem e Desenvolvimento Da Criança e os Campos De Experiências Para embasar a organização do trabalho pedagógico da primeira etapa da Educação Básica, as DCNEI ( Diretrizes Curriculares Nacionais da Educação \textbf{Infantil} ) e a BNCC ( Base Nacional Comum Curricular ) apontam os Direitos de Aprendizagem e Desenvolvimento, como pressupostos importantes da articulação no processo pedagógico no cotidiano da Educação \textbf{Infantil}.
Na organização curricular devem estar contemplados os Direitos de Aprendizagem e Desenvolvimento das \textbf{crianças} previstos para a Educação \textbf{Infantil}, que são : Conviver, \textbf{Brincar}, Participar, Explorar, Expressar e Conhecer -se.
DIREITO DE CONVIVER Conviver com outras \textbf{crianças} e \textbf{adultos}, em \textbf{pequenos} e grandes grupos, utilizando diferentes linguagens, ampliando o conhecimento de si e do outro, o respeito em relação à cultura e às diferenças entre as pessoas.
DIREITO DE \textbf{BRINCAR} \textbf{Brincar} cotidianamente de diversas formas, em diferentes espaços e tempos, com diferentes \textbf{parceiros} ( \textbf{crianças} e \textbf{adultos} ), ampliando e diversificando seu acesso a produções culturais, seus conhecimentos, sua imaginação, sua criatividade, suas experiências emocionais, corporais, \textbf{sensoriais}, expressivas, cognitivas, sociais e \textbf{relacionais}.
DIREITO DE PARTICIPAR Participar ativamente, com \textbf{adultos} e outras \textbf{crianças}, tanto do planejamento da gestão da escola e das atividades propostas pelo educador quanto da realização das atividades da vida cotidiana, tais como a escolha das \textbf{brincadeiras}, dos materiais e dos ambientes, desenvolvendo diferentes linguagens e elaborando conhecimentos, decidindo e se posicionando.
DIREITO DE EXPLORAR Explorar movimentos, \textbf{gestos}, \textbf{sons}, formas, \textbf{texturas}, cores, palavras, emoções, transformações, relacionamentos, histórias, objetos, elementos da natureza, na escola e fora dela, ampliando seus saberes sobre a cultura, em suas diversas modalidades : as artes, a escrita, a ciência e a tecnologia.
DIREITO DE EXPRESSAR Expressar, como sujeito dialógico, criativo e sensível, suas necessidades, emoções, sentimentos, dúvidas, hipóteses, descobertas, opiniões, questionamentos, por meio de diferentes linguagens.
DIREITO DE CONHECER -SE Conhecer -se e construir sua identidade pessoal, social e cultural, constituindo uma imagem positiva de si e de seus grupos de pertencimento, nas diversas experiências de \textbf{cuidados}, interações, \textbf{brincadeiras} e linguagens vivenciadas na \textbf{instituição} escolar e em seu contexto familiar e comunitário.
Os Direitos de Aprendizagem e Desenvolvimento apontados devem ser contemplados considerando os “ Campos de Experiências ” como possibilidades pedagógicas que possibilitem a \textbf{criança} desenvolver ações de descoberta, de forma a assegurar o fazer e o agir da \textbf{criança}.
Assim, cabe ao professor a tarefa de estar \textbf{atento} aos acontecimentos por meio da observação de modo a promover a ampliação da \textbf{curiosidade} da \textbf{criança}, incentivá -la a novas descobertas, contribuindo significativamente, para o processo de aprendizagem da \textbf{criança} a partir da apropriação do seu conhecimento.
Desde \textbf{bebês}, as \textbf{crianças} descobrem diversas formas de comunicação, o que se define como o desenvolvimento das múltiplas linguagens por meio da \textbf{escuta} de vários \textbf{sons} : sua própria voz ou vozes \textbf{ouvidas}, dos jogos sonoros, das músicas, dos seus \textbf{gestos}.
Outros experimentos foram observados em atividades de arte : o símbolo, o desenho, a \textbf{pintura}, \textbf{manipular} materiais e a utilização da mídia favorecem a exploração e a compreensão do mundo.
A arte possibilita a aproximação da realidade, a observação do cotidiano de modo mais consciente para desenvolver ações sobre o meio em que a \textbf{criança} está inserida.
Para o desenvolvimento do cotidiano pedagógico na Educação \textbf{Infantil} podemos considerar outros aspectos que podem contribuir para as práticas de experimentação, a exemplo da exploração de materiais, laboratórios e estúdio de desenho, a exemplo das experiências de Reggio Emília, os objetos do patrimônio, das ruas, das praças, dos jardins, das paisagens, ou ainda da arte mais rebuscada, como as \textbf{pinturas}, esculturas, os museus, arquiteturas, ou mesmo, uma arte percebida ao caminhar no bairro, observar as casas, as cores, os símbolos, etc.
Visto que nenhuma prática vivenciada termina em si mesma, a partir dela produzimos sentidos e um aprendizado constante.
É dessa forma, que se transmite um sentido particular para a compreensão dos Campos de Experiências que devem ser explorados pelas \textbf{crianças} juntas com os professores e demais \textbf{adultos} no ambiente educativo da Educação \textbf{Infantil}.
Assim, ao pensar em Campos de Experiências não podemos relacionar ao modelo de currículo estruturado por meio das divisões de áreas do conhecimento.
Não se trata de um simples olhar isolado por disciplina com uma organização fragmentada da realidade, mas compreender que os Campos de Experiências constituem um amplo contexto que envolve toda organização do trabalho pedagógico, ( tempo, espaço, escolha dos materiais, trabalho em grupo e o \textbf{acompanhamento} da aprendizagem das \textbf{crianças}, como uma “ Ecologia educativa ”, o que implica abranger também os instrumentos e os artefatos culturais, as imagens e as palavras.
O termo “ Campo de Experiências ” utilizado na organização do currículo surgiu da necessidade de centralizar as ações das \textbf{crianças}, ou seja, torná -las protagonistas do processo educativo.
Considerar “ experiência ” como atividades contínuas e participações ativas entre as \textbf{crianças} valoriza as dimensões de suas ações com o universo dos patrimônios da humanidade em relação a complexidade e a transversalidade.
Os Campos de Experiências procuram vincular os direitos das \textbf{crianças} aos conhecimentos já sistematizados.
Colocam as interações e as \textbf{brincadeiras} no centro do processo educativo das quais emergem as ações das \textbf{crianças}, as observações, as investigações, os posicionamentos e as significações.
Cada Campo de Experiências oferece às \textbf{crianças} oportunidades de interagir com pessoas, objetos, situações, atribuindo um sentido pessoal a essas interações.
Tendo a experiência um sentido específico para cada \textbf{criança}, ao planejar um contexto educativo, o professor cria formas de \textbf{registros} para acompanhar o desenvolvimento da aprendizagem da \textbf{criança}, através do qual fornecer-lhe-á pistas para a continuidade do trabalho pedagógico.
Nesse sentido, os Campos de Experiências não focam apenas a \textbf{criança}, mas as relações que ocorrem entre seus \textbf{pares}, o professor, os familiares, a comunidade, os saberes, as linguagens, o conhecimento e o mundo, promovendo a socialização das \textbf{crianças} como um importante processo de aprendizagem e desenvolvimento.
Os Campos de Experiências em que se estrutura a BNCC, são : - O eu, o outro eu nós ; - Corpo, \textbf{gestos} e movimentos ; - \textbf{Escuta}, fala, linguagem e pensamento ; - Traços, \textbf{sons}, cores e formas ; - Espaços, tempos, quantidades, relações e transformações.
Os Campos de Experiências são abordados de forma articulada e mediados pelas perspectivas das interações e das \textbf{brincadeiras} como eixo de organização da aprendizagem e desenvolvimento.
É na relação com o outro que o sujeito se constitui como ser individual nas práticas sociais das quais participa e na cultura em que está inserido.
Nesse processo, o sujeito vai se transformando e mudando o mundo à sua \textbf{volta} na relação com o Campo “ O eu, o outro e o nós ”.
Nessa perspectiva nos questionamos, “ Que mundo estamos apresentando às \textbf{crianças}?
É preciso ter em mente que as \textbf{crianças} imitam \textbf{gestos}, expressões, falas, atitudes e modos de ser e de conviver dos \textbf{adultos} e de outras \textbf{crianças} com as quais convivem.
É por essa razão que no trabalho pedagógico as relações pessoais são tão importantes.
São elas que irão possibilitar às \textbf{crianças} uma relação saudável com o mundo.
A abordagem do Campo de Experiências “ Corpo, \textbf{gestos} e movimentos ” se refere à linguagem corporal.
O corpo não se limita às características físicas e biológicas, mas expressam aspectos que revelam o grupo social a que pertencemos, nossa singularidade, nossa identidade pessoal e social.
Com o corpo, a \textbf{criança}, desde \textbf{bebê}, por meio dos seus sentidos e movimentos explora o mundo e constrói conhecimentos sobre si, o outro e o universo social e cultural.
Nesse processo, as \textbf{crianças}, ao \textbf{brincar} com seu corpo, desenvolvem múltiplas linguagens através da música, da dança, do teatro, das \textbf{brincadeiras} de faz-de-conta.
Nas práticas pedagógicas das \textbf{instituições} de Educação \textbf{Infantil}, o corpo é o elemento central.
O movimento corporal explora as possibilidades necessárias para a \textbf{criança} identificar suas potencialidades e seus limites, o que permite uma atuação mais autônoma da \textbf{criança} no seu meio.
O Campo de Experiências “ \textbf{Escuta}, fala, pensamento e imaginação ” revela o choro do \textbf{bebê} como a primeira linguagem oral do sujeito.
Além do choro, o \textbf{bebê} se utiliza de outras formas de comunicação, como os olhares e os \textbf{gestos} e vai se relacionando com uma infinidade de \textbf{sons} \textbf{emitidos} por outros sujeitos de sua cultura, ampliando o seu repertório vocal.
Posteriormente, a partir dos balbucios, as \textbf{crianças} vão se apropriando das palavras ao tentarem pronunciá -las com auxílio dos objetos.
Nessa associação entre os \textbf{sons} e as palavras surgem a internalização e a constituição do pensamento verbal, ou seja, a \textbf{criança} passa a expressar suas ações por meio das palavras, explorando -as como instrumento essencial de pensamento e comunicação.
Nessa dinâmica as unidades de Educação \textbf{Infantil} devem envolver as \textbf{crianças} em atividades da cultura oral, como a \textbf{escuta} de histórias, \textbf{roda} de conversas, descrições, narrativas, permitindo -as explorar as múltiplas linguagens.
Assim, como as \textbf{crianças} quando inseridas em contextos comunicativos adquirem a linguagem oral, ao vivenciarem situações significativas de uso de leitura e escrita, iniciam o processo de aprendizagem dessas linguagens.
Para Salles e Farias ( 2012, p.136 ), > a Educação \textbf{Infantil} tem o importante papel de possibilitar o acesso das \textbf{crianças} a cultura letrada, a partir da vivência de experiências com diversos suportes e gêneros textuais, em práticas sociais reais em que o uso desses textos se torne necessário.
O processo de aprendizagem dessa linguagem pelas \textbf{crianças} nessa etapa da Educação Básica se amplia, na medida em que são trabalhados, de modo intencional, os processos de produção e de leitura de textos.
O Campo de Experiências “ Traços, \textbf{sons}, cores e formas ” envolve experiências relacionadas ao fazer artístico, à apreciação estética, à reflexão e à apropriação de conhecimentos das artes.
Essas experiências, ao serem desenvolvidas na Educação \textbf{Infantil}, abrangem várias modalidades que se expressam por meio de múltiplas linguagens, como as artes visuais, a dança, a música e o teatro.
A aquisição de diferentes manifestações artísticas e culturais é fundamental para a formação humana por possibilitar a \textbf{criança} expandir suas experiências de vida e recriar um repertório de práticas e conceitos singulares.
Por estar imersa em um mundo construído de fenômenos naturais e socioculturais, a \textbf{criança}, desde \textbf{bebê}, percebe o espaço que ocupa e procura se situar no tempo.
Revela \textbf{curiosidades} sobre o mundo físico, em relação ao seu próprio corpo, aos fenômenos atmosféricos e as transformações da natureza.
A \textbf{criança} também faz indagações a respeito do mundo sociocultural envolvendo as relações familiares e entre outras pessoas com as quais convive.
Nessas práticas, ela ainda vivencia conceitos matemáticos referentes às relações quantitativas, conceito de número, medidas, grandezas, formas e relações espaço/tempo.
Essas abordagens constituem o Campo de Experiências “ Espaços, tempos, quantidades, relações e transformações ” que abrange tanto os conhecimentos do cotidiano, como também os conhecimentos historicamente acumulados pela humanidade.
Nesse sentido, cabe à \textbf{instituição} de Educação \textbf{Infantil} proporcionar práticas que possibilitem à \textbf{criança} buscar respostas às suas indagações e \textbf{curiosidades} por meio de experiências que lhe permitam observar, \textbf{manipular} objetos, levantar hipóteses, consultar fontes de informações, investigar e explorar o mundo ao seu redor.
Dessa forma, a \textbf{criança} terá oportunidade para ampliar seus conhecimentos do mundo físico e sociocultural e utilizá -los tanto em seu \textbf{dia} a \textbf{dia} como em seus estudos posteriores.

% matched lemmas: acompanhamento, adulto, atento, bebê, brincadeira, brincar, criança, cuidado, curiosidade, dia, emitir, escuta, gesto, infantil, instituição, manipular, ouvir, par, parceiro, pequeno, pintura, registro, relacional, roda, sensorial, som, textura, volta
\end{textsample}
